\documentclass[tikz,border=4pt]{standalone}
\usepackage{ctex}
\usepackage{amsmath,amssymb}
\usetikzlibrary{positioning}
\begin{document}
\begin{tikzpicture}[
  hdr/.style={rectangle, draw=blue!70, fill=blue!15, font=\small\bfseries,
              minimum width=5.5cm, minimum height=0.7cm, text centered},
  row/.style={rectangle, draw=gray!40, fill=white,
              minimum width=5.5cm, minimum height=0.7cm, text centered, font=\footnotesize},
  altrow/.style={rectangle, draw=gray!40, fill=gray!6,
              minimum width=5.5cm, minimum height=0.7cm, text centered, font=\footnotesize},
  hdr2/.style={rectangle, draw=orange!70, fill=orange!15, font=\small\bfseries,
              minimum width=4.8cm, minimum height=0.7cm, text centered},
  row2/.style={rectangle, draw=gray!40, fill=white,
              minimum width=4.8cm, minimum height=0.7cm, text centered, font=\footnotesize},
  altrow2/.style={rectangle, draw=gray!40, fill=gray!6,
              minimum width=4.8cm, minimum height=0.7cm, text centered, font=\footnotesize},
]

% ===== Left table: Derivative formulas =====
\node[hdr] (h1) at (0,0) {矩阵求导速查（分母布局）};
\node[row,  below=0pt of h1] (r1) {$\nabla_x (a^\top x) = a$};
\node[altrow, below=0pt of r1] (r2) {$\nabla_x (x^\top A x) = (A+A^\top)x$};
\node[row,  below=0pt of r2] (r3) {$\nabla_x (x^\top A x) = 2Ax$（$A$ 对称）};
\node[altrow, below=0pt of r3] (r4) {$\dfrac{\partial\,\mathrm{tr}(AB)}{\partial A}=B^\top$};
\node[row,  below=0pt of r4] (r5) {$\dfrac{\partial\,\ln|A|}{\partial A}=A^{-T}$};
\node[altrow, below=0pt of r5] (r6) {$d(A^{-1})=-A^{-1}dA\,A^{-1}$};
\node[row,  below=0pt of r6] (r7) {$\dfrac{\partial(Ax)}{\partial x}=A$};
\node[altrow, below=0pt of r7] (r8) {Softmax: $J=\mathrm{diag}(p)-pp^\top$};

% ===== Right table: Layout conventions =====
\node[hdr2] (h2) at (7.5,0) {布局约定对比};
\node[row2,  below=0pt of h2] (s1) {\textbf{分母布局（本书约定）}};
\node[altrow2, below=0pt of s1] (s2) {$\nabla_x f\in\mathbb{R}^n$ 为列向量};
\node[row2,  below=0pt of s2] (s3) {$J_{ij}=\dfrac{\partial y_i}{\partial x_j}$};
\node[altrow2, below=0pt of s3] (s4) {\textbf{分子布局（部分工程约定）}};
\node[row2,  below=0pt of s4] (s5) {$\nabla_x f\in\mathbb{R}^{1\times n}$ 为行向量};
\node[altrow2, below=0pt of s5] (s6) {两者转置关系，混用必出错};
\node[row2,  below=0pt of s6] (s7) {\textbf{链式法则（分母布局）}};
\node[altrow2, below=0pt of s7] (s8) {$\nabla_x f=J^\top_{y/x}\nabla_y f$};

% Tip box
\node[rectangle, draw=red!60, fill=red!6, font=\footnotesize, align=left,
      minimum width=12.5cm, minimum height=0.9cm, rounded corners]
  at (3.7, -6.4)
  {注意：$A$ 对称时 $\nabla_x(x^\top Ax)=2Ax$；若 $A$ 不对称则为 $(A+A^\top)x$。
   Hessian $=\nabla^2_x(x^\top Ax)=A+A^\top$（对称情形 $=2A$）。};

\node[font=\normalsize\bfseries] at (3.7, 0.9) {矩阵微积分公式速查表};
\end{tikzpicture}
\end{document}
